%% AUTOMATISCH ERZEUGT (make_ki_anhang.py) -- Prompts und Inputs der KI-Nutzung zu Einlage 1
\chapter{KI-Nutzung zu Einlage 1: Prompts und Inputs}
\label{anh:ki-arbeitsplan}

Beschreibung und Reflexion der Einlage Arbeitsplan (Abschnitt~\ref{sec:arbeitsplan}) sind von einem
KI-System formuliert worden \cite{Manevski2026KIArbeitsplan}. Verwendetes System: Claude Opus 5.5
(Anthropic) in Claude Code, 28.\,10.\,2026. Weitere Inputs: der gemeinsame Arbeitsplan
(Arbeitsplan\_Wackelstein.xlsx), die Kursunterlagen im Projektordner und die Arbeitsplan-Abgabe von
Janosch Lustig, die nur als Vergleich für Umfang und Stil diente. Im Folgenden stehen alle meine Prompts
in der Reihenfolge, in der ich sie eingegeben habe, und danach der von der KI erstellte und von mir
ausgefüllte Fragebogen.

\section*{Prompts}
{\small

\par\medskip\noindent\textbf{Prompt 1:}\\
{\ttfamily\raggedright
schau dir Moritz Beispielportfolio an (C:\textbackslash{}Users\textbackslash{}KIKO\textbackslash{}Documents\textbackslash{}!UNI\textbackslash{}Projekt-Praktikum\textbackslash{}Formulare für Abgaben\textbackslash{}Moritz Beispiel)\\
\smallskip
passe alle datei darin an sodass es für mich zugeschnitten ist und alle informatione beinhaltet die im projekt zu finden sind\\
speichere meine angepasste Kopie im ''Formulare für Abgaben'' Ordner als Kristijan Protfolio\\
}
\par\medskip\noindent\textbf{Prompt 2:}\\
{\ttfamily\raggedright
Nicht jetzt\\
}
\par\medskip\noindent\textbf{Prompt 3:}\\
{\ttfamily\raggedright
ok jetzt eine neue aufgabe\\
findest du in dem projekt eine datei die beschreibt wie die Reflexion des Arbeitsplans aussehen soll (inhalt, länge, etc)?\\
}
\par\medskip\noindent\textbf{Prompt 4:}\\
{\ttfamily\raggedright
{[}Eingefügter Text der Moodle-Aufgabe:{]}\\
Einlage: Arbeitsplan – Aufgabe\\
Termin: 28. Oktober 2026, 23:59 Uhr\\
\smallskip
1. Erstellen Sie gemeinsam mit Ihrer Partnerin / Ihrem Partner einen Arbeitsplan für Ihr Projekt. Besprechen Sie den Plan mit den Betreuern. Als Ideengeber können Sie den Arbeitsplan zum Mittagessen kochen nehmen. Es ist sinnvoll, den Plan im Gantt-Format zu erstellen. Neben spezialisierter Software können Sie auch Excel oder Libre Office Calc verwenden.\\
2. Schreiben Sie eine Reflexion über den Arbeitsplan. Eine Reflexion ist keine Nacherzählung. Verwenden Sie die zur Verfügung gestellten Template oder strukturieren Sie Ihre Abgabe danach.\\
3. Für diese Einlage verwenden alle Gruppenmitglieder den gleichen gemeinsamen Arbeitsplan, schreiben jedoch die Reflexionen einzeln und unabhängig voneinander.\\
4. Laden Sie den Arbeitsplan und die Reflexion hoch.\\
   * Akzeptierte Formate für alle Abgaben: Adobe PDF (Microsoft DOCX bitte zuert in PDF speichern). Wandeln Sie andere Formate in diese Formate um.\\
   * Umfang Beschreibung der Einlage: Mindestens: 1/4-Seite, Maximal 1/2 Seite\\
   * Umfang Reflexion: Mindestens 1 Seite, maximal 2 Seiten!\\
   * Alle Längenangaben sind ohne Titel, Abbildungen, Abbildungslegenden, Tabellen und Tabellenlegenden (oder Ähnliches) zu verstehen.\\
5. Vergessen Sie nicht nach dem Hochladen den Abgabeknopf zu drücken.\\
\smallskip
Dazugehörige Dateien finden Sie in Unterlagen zum Projektpraktikum, Einlage während Semester Template und Unterlagen zum Projektpraktikum, Arbeitsplan\\
}
\par\medskip\noindent\textbf{Prompt 5:}\\
{\ttfamily\raggedright
ok für den arbeitsplan\\
\smallskip
1. Einlage:\\
   1. was zur hölle ist eine ''Einlage''\\
   2. Stelle mir einen Fragebogen den ich beantworten soll damit du daraus die Einlage schreiben kannst\\
2. Reflexion:\\
   1. was ist mit ''Reflexion'' gemeint\\
   2. Stelle mir einen riesen fragebogen den ich beantworten soll damit du damit die Reflexion schreiben kannst\\
}
\par\medskip\noindent\textbf{Prompt 6:}\\
{\ttfamily\raggedright
{[}Anhang: pp\_reflexion\_arbeitsplan (1).pdf – Arbeitsplan-Abgabe von Janosch Lustig{]}\\
die .md file habe ich nun ausgefüllt, lesen sie nochmals durch und mache damit den text.\\
\smallskip
hier ist noch die in der .md angesprochene version von janosch\\
\smallskip
übertreib es nicht mit der länge\\
}
\par\medskip\noindent\textbf{Prompt 7:}\\
{\ttfamily\raggedright
Mit KI-Zitat schreiben\\
}
}

\section*{Fragebogen mit meinen Antworten}
{\small\raggedright

\par\noindent Antworte direkt unter jede Frage, \textbf{Stichpunkte reichen}. Ehrlich und konkret ist besser als schön:\par
\par\noindent „Ich hab die E-Werkstatt-Wartezeit unterschätzt, weil …“ ist Gold, „Planung ist wichtig“ ist wertlos.\par
\par\noindent Fragen, zu denen dir nichts einfällt: leer lassen oder „weiß nicht“. Nichts wird erfunden.\par
\subsection*{Teil 1 – Beschreibung der Einlage (1/4 bis 1/2 Seite)}
\par\noindent Ziel: Ein Leser soll wissen, \emph{was} der Arbeitsplan ist, \emph{woher} er kommt und \emph{was drinsteht} –\par
\par\noindent ohne ihn nachzuerzählen.\par
\par\noindent \textbf{1.1} Womit habt ihr den Plan gebaut? (Excel? Version? Mit Vorlage oder komplett selbst?)\par
\par\noindent\emph{Antwort:} Wir haben den Plan in Excel Gebaut. Haben uns an einer Baumstruktur pro parallel laufendem thema ein beispiel genommen.\par\smallskip
\par\noindent \textbf{1.2} Wann genau habt ihr ihn erstellt? (erste Version an welchem Tag, letzte Änderung? Der Plan sagt „Stand 22.10.2026“ – stimmt das noch?)\par
\par\noindent\emph{Antwort:} Der Plan wurde seit dem 22.10.2026 nicht verändert und wurde an genau diesem tag fertiggestellt und mit Prof. Marti durchgesprochen\par\smallskip
\par\noindent \textbf{1.3} Wer hat was daran gemacht? (Struktur, Aufgabenliste, Formatierung/Formeln, Termine-Blatt …)\par
\par\noindent\emph{Antwort:} Komplette zusammenarbeit in jedem bereich des plans\par\smallskip
\par\noindent \textbf{1.4} Wie lange habt ihr ungefähr daran gearbeitet (Stunden, wie viele Sitzungen)?\par
\par\noindent\emph{Antwort:} nicht getrackt\par\smallskip
\par\noindent \textbf{1.5} Was habt ihr als Vorlage / Ideengeber benutzt? (Beispiel „Mittagessen kochen“ aus Moodle? Petros' Plan – was genau habt ihr davon übernommen: Struktur, Termine, Idee mit Python/CSV?)\par
\par\noindent\emph{Antwort:} Mittagessen kochen als beispiel für die detalierten zweige des baums\par\smallskip
\par\noindent \textbf{1.6} Habt ihr den Plan mit den Betreuern besprochen? Wenn ja: wann, mit wem (Name), und was haben sie gesagt?\par
\par\noindent\emph{Antwort:} ja wie gesagt mit prof. marti und dieser hat ein paar fragen gestellt die wir im beantworten konnten. dannach hat er abgenickt\par\smallskip
\par\noindent \textbf{1.7} Was ist für dich die wichtigste Aussage des Plans in einem Satz? (z. B. „Kritischer Pfad ist der Sensor“, „Messungen erst ab Dezember möglich“ …)\par
\par\noindent\emph{Antwort:} für mich ist die wichtigeste aussage des plans die super leicht zu erkennenden aufgaben die parallel laufen und zusätzlich auch noch zusammenhängen damit man möglichst effizient ist\par\smallskip
\par\noindent \textbf{1.8} Welche Phasen / Meilensteine sind aus deiner Sicht die entscheidenden?\par
\par\noindent\emph{Antwort:} erster meilenstein wäre der fertige prototyp (das er funktioniert und man damit anfangen kann zu messen), zweiter meilenstein wäre die funktionierende simulation sodass man seine ergebnisse direkt validieren/vergleichen kann, dritter meilenstein ist dann eine schöne saubere messreihe womit man ordentlich arbeiten kann\par\smallskip
\par\noindent \textbf{1.9} Gibt es Besonderheiten am Plan, die erwähnt werden sollten? (Farbcodierung nach Person, automatische Balken, Statusspalte, Termine-Blatt mit Quellen, Praktikumspause …)\par
\par\noindent\emph{Antwort:} alles ist nach person aufgeteilt, die statusspalte ist super zum nachverfolgen was der andere macht, die prozentangabe des themenblocks ist praktisch um sich zu motivieren oder nichts zu vergessen und alle abgaben sind eingetragen\par\smallskip
\par\noindent \textbf{1.10} Ist der Plan als „lebendes Dokument“ gedacht? Wie oft wollt ihr ihn aktualisieren?\par
\par\noindent\emph{Antwort:} der plan darf nicht verändert werden. er wird von uns nur dazu genutzt um unseren fortschritt zu tracken.\par\smallskip
\subsection*{Teil 2 – Reflexion (1 bis 2 Seiten)}
\par\noindent Ziel: Nicht \emph{was} im Plan steht, sondern \emph{was das Erstellen mit dir gemacht hat} – was du gelernt,\par
\par\noindent gemerkt, falsch eingeschätzt hast und was du daraus folgerst.\par
\par\medskip\noindent\textbf{A – Wie lief das Erstellen?}\par
\par\noindent \textbf{A1} Wie seid ihr vorgegangen? (erst Liste, dann Zeiten? Rückwärts von den Abgaben geplant? Brainstorming?)\par
\par\noindent\emph{Antwort:} wir konnten anfangs nicht in worte fassen was wir uns vorgestellt haben wie der plan auszusehen hat. dann haben wir gebrainstormed und ''prototypen'' gezeichnet.\par\smallskip
\par\noindent \textbf{A2} Was war der erste Entwurf, und wie sehr hat sich der Plan bis zur jetzigen Version verändert? Was wurde umgeworfen und warum?\par
\par\noindent\emph{Antwort:} der plan ist nicht direkt in excel entstanden sondern im kopf dann mehrere itterationen auf dem blatt und dann erst als grober ansatz in excel\par\smallskip
\par\noindent \textbf{A3} Wie habt ihr entschieden, wer welche Aufgabe übernimmt? Nach Interesse, Können, Zeit?\par
\par\noindent\emph{Antwort:} unsere themenschwerpunkte waren uns klar. daher haben wir uns zu jedem schwerpunkt überlegt welche aufgaben alle zu erledigen sind. diese aufgaben haben wir dann alle runtergeschrieben. und untereinander aufgeteilt. dannach haben wir überlegt welche aufgaben gleichzeitig zusammenhängen und in welchem zeitabschnitt des PPs sie zu erledigen sind damit es am sinnvollsten ist. so haben wir dann die aufgaben gruppiert und aufgeteilt wann diese zu erledigen sind. jedoch können wir ja nicht jeder aufgabe gleich viel zeit geben, also haben wir noch disskutiert welche aufgabe wie viel zeit in anspruch nehmen wird, wobei die GP erfahrung der letzten 2 Jahre geholfen hat.\par\smallskip
\par\noindent \textbf{A4} Gab es Diskussionen oder Meinungsverschiedenheiten mit Janosch? Worüber? Wie gelöst?\par
\par\noindent\emph{Antwort:} es gab direkt am anfang meinungsverschiedenheiten die wir durch argumente und vielleicht vergessene themen alles schön lösen konnten sodass am ende beide glücklich sind und das beste aus beiden welten zusammenfindet\par\smallskip
\par\noindent \textbf{A5} Wie lief die Arbeit mit Excel selbst? (Formeln, bedingte Formatierung – leicht, nervig, neu für dich?)\par
\par\noindent\emph{Antwort:} frage ist todes unwichtig\par\smallskip
\par\medskip\noindent\textbf{B – Bezug zum Projektthema (Wackelstein)}\par
\par\noindent \textbf{B1} Was wusstest du über den Wackelstein / das Projekt \emph{bevor} ihr den Plan gemacht habt, und was danach?\par
\par\noindent\emph{Antwort:} ich wusste davor nur das er existiert und grob wie seine funktionalität ist. nach dem erstellen des plans ist mir die aufteilung des verständnisses klarer geworden. wie man eine funktion des steins betrachten sollte damit man sie am ende intuitiver oder besser versteht und was sich alles in der theorie versteht an das man nicht sofort nur beim anschauen des steins erkennt, was jedoch klar wird wenn man sich ein ziel für den versuch überlegt und dann die aufgaben ausdenken muss die einen zum ziel führen.\par\smallskip
\par\noindent \textbf{B2} Welche Arbeitsschritte sind dir erst beim Planen bewusst geworden? (z. B. Glasplatte testen, Reibung, Nivellieren, reproduzierbarer Anstoß, Sensor kalibrieren)\par
\par\noindent\emph{Antwort:} reibung, reproduzierbarer anstoß, sonsorkalibrierung nie als konkreten schritt angesehen welcher aber nach dem daran denken unfassbar wichtig und praktisch ist\par\smallskip
\par\noindent \textbf{B3} Wo hat das Planen dir gezeigt, dass das Projekt aufwendiger / einfacher ist als gedacht?\par
\par\noindent\emph{Antwort:} bei der theorie und der vorbereitung (das 3D model muss für so einen niechigen effekt sozusagen PERFEKT sein)\par\smallskip
\par\noindent \textbf{B4} Welche Teile des Projekts hängen voneinander ab, und hast du das vorher gesehen? (Sensor $\rightarrow$ Messkörper $\rightarrow$ Messungen; Trägheitstensor $\rightarrow$ Simulation)\par
\par\noindent\emph{Antwort:} die messungen hängen von der qualität des körpers ab, was klar ist. jedoch war es vorerst unklar das die kleinste veränderung der startbedingungen eine so wichtige rolle spielen wird.\par\smallskip
\par\medskip\noindent\textbf{C – Bezug zu deinem Schwerpunkt (Bewegungsgleichungen + Simulation der Spin-Umkehr)}\par
\par\noindent \textbf{C1} Ist dieser Schwerpunkt-Titel richtig? Wenn nicht: wie heißt dein Schwerpunkt genau, und was ist der von Janosch?\par
\par\noindent\emph{Antwort:} das ist meiner ja\par\smallskip
\par\noindent \textbf{C2} Wo taucht dein Schwerpunkt im Plan auf (welche Zeilen), und wie viel Zeit hast du dafür eingeplant? Realistisch?\par
\par\noindent\emph{Antwort:} das sollst du rausfinden (claude), das steht im plan\par\smallskip
\par\noindent \textbf{C3} Was brauchst du von Janosch, damit du deinen Teil machen kannst (z. B. Trägheitstensor aus der STL), und was braucht er von dir?\par
\par\noindent\emph{Antwort:} alle wichtigen größen des körpers und seine eigenschaften.\par\smallskip
\par\noindent \textbf{C4} Hat dir der Plan geholfen zu verstehen, was dein Schwerpunkt konkret bedeutet? Was war vorher unklar?\par
\par\noindent\emph{Antwort:} ja hat er da ich die chance hatte mich genauer damit zu befassen wie ich meine aufgaben einzuteilen habe und womit sie direkt oder indirekt zusammenhängen.\par\smallskip
\par\noindent \textbf{C5} Warum hast du (bzw. habt ihr) gerade dir die Aufgaben Bewegungsgleichungen, Simulation, Firmware, Kalibrierung, Datenpipeline gegeben? Was reizt dich daran?\par
\par\noindent\emph{Antwort:} mich reizt die programmierungm, theorie und simulaton da man dabei ein tieferes verständnis bekommt da man sich um die kleinstens details kümmern muss an die man sonst, wenn man zb eine simulation nicht als ziel hätte, nie denken würde und die dann am ende im verständnis fehlen würden\par\smallskip
\par\medskip\noindent\textbf{D – Was hast du gelernt?}\par
\par\noindent \textbf{D1} Was hast du über Projektplanung allgemein gelernt? (Aufgaben zerlegen, Puffer, Parallelität, externe Abhängigkeiten …)\par
\par\noindent\emph{Antwort:} das puffer sehr wichtig ist wenn man einen gegebenen zeitraum und etwas unbekanntes vor sich hat, es kann ja schlieslich alles schief gehen. gelernt habe ich auch das es sehr wichtig ist zu itterieren, wenn man an einem schon bekanntem thema nach einiger zeit und damit mit einem andern winkel arbeitet erkannt man was man vergessen hat oder hinterfragt die richtigkeit und wichtigkeit des schon erarbeitetem.\par\smallskip
\par\noindent \textbf{D2} Was hast du über dich selbst gelernt? (Neigst du zum Unterschätzen? Planst du gern / ungern? Arbeitest du lieber spontan?)\par
\par\noindent\emph{Antwort:} ich habe über mich gelern das ich eher dazu neige die zeit die man für etwas hat zu unterschätzen. dessen bin ich mir jedoch schon bewusst wordurch ich dafür kompensieren kann, so wie in diese plan zu sehen ist.\par\smallskip
\par\noindent \textbf{D3} Hattest du vorher schon mal so einen Plan gemacht (Schule, Job, andere Praktika)? Unterschied?\par
\par\noindent\emph{Antwort:} keinen solch detalierten plan\par\smallskip
\par\noindent \textbf{D4} Welche Vorlesung / welches Modul hätte dir dabei geholfen oder hat geholfen?\par
\par\noindent\emph{Antwort:} dabei hat die erfahrung aus dem GP (Grundpraktikum) geholfen da ich dort die erfahrung gemacht habe das aufgaben bei denen man denkt das sie schnell erledigt sind eigentlich nie schnell erledigt sind, wodurch ich in diesem plan (hoffentlich) nicht unterschätzt habe\par\smallskip
\par\medskip\noindent\textbf{E – Schwierigkeiten}\par
\par\noindent \textbf{E1} Was war am schwersten? (Dauer schätzen, Reihenfolge, wissen was man noch nicht weiß, Excel, Absprache …)\par
\par\noindent\emph{Antwort:} ich würde nicht sagen das mir etwas sehr schwer fiel. am anfang war der gedanke einen solch detalierten plan zu erstellen einschüchternd, hat aber nach dem beginnen nur positives gebracht wodurch am ende eher zufriedenheit als frust oder schwierigkeit zu emfinden war.\par\smallskip
\par\noindent \textbf{E2} Welche Dauern konntest du gar nicht einschätzen, und wie habt ihr dann geschätzt?\par
\par\noindent\emph{Antwort:} da es schwer ist mit dem 3D druck und teilliefenrungen zu planen haben wir dafür ordentlich puffer eingebaut damit es an solchen banalen dingen nicht scheitert.\par\smallskip
\par\noindent \textbf{E3} Wo hängt ihr von anderen ab (E-Werkstatt, 3D-Druck, Betreuer, Material)? Wie seid ihr damit umgegangen?\par
\par\noindent\emph{Antwort:} von unserem betreuer da er unseren Stein 3D druckt und wir hoffen müssen das er dauerhaft erreichbar ist falls es änderungen oder sowas gibt\par\smallskip
\par\noindent \textbf{E4} Gab es Aufgaben, bei denen du beim Planen nicht wusstest, was sie eigentlich beinhalten?\par
\par\noindent\emph{Antwort:} nicht wirlich, und auch wenn das am anfang so gewesen sein sollte ist es noch WÄHREND dem planen klar geworden.\par\smallskip
\par\medskip\noindent\textbf{F – Feedback von anderen}\par
\par\noindent \textbf{F1} Was hat Janosch zum Plan bzw. zu deinen Ideen gesagt?\par
\par\noindent\emph{Antwort:} es gab kritik sowe positives, an genaue wortlaute kann ich mich nicht errinnern. am ende waren wir beide zufrieden, was am wichtigsten ist\par\smallskip
\par\noindent \textbf{F2} Was haben die Betreuer gesagt? (Name, Datum, Inhalt – wird zitiert!)\par
\par\noindent\emph{Antwort:} wir haben wie gesagt mit marit gesprochen, dieser hat den plan in der gezeigten version abgenickt nachdem er fragen gestellt hat um den plan genauer zu verstehen (kein zitat, das gespräch war mündlich am Donnerstag den 22.10.2026)\par\smallskip
\par\noindent \textbf{F3} Hast du den Plan Kommilitonen gezeigt (Petros oder andere)? Was kam als Kommentar? (Name, Datum – wird zitiert!)\par
\par\noindent\emph{Antwort:} nicht wirklich außer das der plan anfangs hässlich war\par\smallskip
\par\noindent \textbf{F4} Was hast du aus dem Feedback geändert – oder bewusst nicht geändert, und warum?\par
\par\noindent\emph{Antwort:} wir haben den plan nur schöner gemacht, an der sturktur jedoch nichts verändert\par\smallskip
\par\medskip\noindent\textbf{G – Fehler und Schwachstellen im Plan}\par
\par\noindent Ich habe im Excel-Plan folgende Auffälligkeiten gefunden. Bitte kommentieren: echter Fehler, Absicht, oder schon behoben?\par
\par\noindent \textbf{G1} CAD des Sensorfachs (2.2.2.1) ist am 09.11., aber „Drucken + Sensor einbauen“ (2.2.2.2) beginnt schon am 05.11.\par
\par\noindent\emph{Antwort:} kein bck mir das nochmal anzuschauen, darüber haben wir jedoch geredet und es macht sinn\par\smallskip
\par\noindent \textbf{G2} Messreihen 1–3 laufen bis 21.01., die „Endauswertung inkl. Messreihe 3“ endet aber am 19.01., und das Messergebnis ist am 20.01. abzugeben.\par
\par\noindent\emph{Antwort:} das ist puffer oder für nachmessungen\par\smallskip
\par\noindent \textbf{G3} Folien erstellen (ab 14.01.) überschneidet sich mit den letzten Messungen.\par
\par\noindent\emph{Antwort:} das passt, messungen dauern nicht lange\par\smallskip
\par\noindent \textbf{G4} Paper durcharbeiten nur 05.–10.11. vor der Abgabe am 11.11. – zu knapp?\par
\par\noindent\emph{Antwort:} nope das passt auch\par\smallskip
\par\noindent \textbf{G5} „Reproduzierbarer Anstoß“ (3.3) und „E-Werkstatt-Bearbeitung“ haben „offen“ als Verantwortlichen – Absicht?\par
\par\noindent\emph{Antwort:} ja das ist richtig\par\smallskip
\par\noindent \textbf{G6} Siehst du selbst weitere Fehler, Lücken oder zu optimistische Zeiten? Fehlt eine Aufgabe ganz (z. B. Theorie-Vergleich vorbereiten, Backup der Daten, Akku laden …)?\par
\par\noindent\emph{Antwort:} nein\par\smallskip
\par\noindent \textbf{G7} Gibt es physikalische Annahmen im Plan, die vielleicht falsch sind? (z. B. dass die Glasplatte Rollen ohne Gleiten erlaubt, dass ein definiertes Anfangs-$\omega$ überhaupt reproduzierbar ist)\par
\par\noindent\emph{Antwort:} das ist noch rauszufinden\par\smallskip
\par\medskip\noindent\textbf{H – Überraschungen}\par
\par\noindent \textbf{H1} Gab es eine überraschende Erkenntnis beim Planen? (über das Projekt, den Zeitbedarf, die Zusammenarbeit, dich selbst)\par
\par\noindent\emph{Antwort:} nicht wirklich. diese zusammenarbeit findet seit 2 Jahren in der konstelation statt.\par\smallskip
\par\noindent \textbf{H2} Hat dich irgendwas am Umfang des Projekts erschreckt oder beruhigt?\par
\par\noindent\emph{Antwort:} ich dachte erst das ich wenig spaß an dem projekt haben werde, mittleiweile glaube ich das ich ganz schön viel spaß haben werde nachdem ich mir meine aufgaben deutlich ausgearbeitet habe\par\smallskip
\par\medskip\noindent\textbf{I – Konsequenzen und Ausblick}\par
\par\noindent \textbf{I1} Was würdest du beim nächsten Arbeitsplan anders machen?\par
\par\noindent\emph{Antwort:} das entscheide ich nachdem ich versucht habe diesen zu erfüllen und mir meine kritik an meinen eigenen plan gemerkt habe\par\smallskip
\par\noindent \textbf{I2} Was änderst du konkret an deinem Vorgehen in den nächsten Wochen wegen des Plans?\par
\par\noindent\emph{Antwort:} nichts\par\smallskip
\par\noindent \textbf{I3} Wo erwartest du, dass der Plan nicht halten wird? Was ist dein Plan B?\par
\par\noindent\emph{Antwort:} keine ahnunng\par\smallskip
\par\noindent \textbf{I4} Wie wollt ihr kontrollieren, ob ihr im Zeitplan seid (wöchentlicher Check, Status-Spalte …)?\par
\par\noindent\emph{Antwort:} ja wir stehen in ständiger kommunikation und behalten die aktuellen aufgaben im blick\par\smallskip
\par\noindent \textbf{I5} Welche neue Frage ist für dich durch das Planen aufgetaucht?\par
\par\noindent\emph{Antwort:} keine\par\smallskip
\par\medskip\noindent\textbf{J – Stil}\par
\par\noindent \textbf{J1} Ich-Form, eher sachlich oder eher persönlich?\par
\par\noindent\emph{Antwort:} ich form, nicht zu sachlich aber eher sachlich mit persönlichkeit reingestreut\par\smallskip
\par\noindent \textbf{J2} Gibt es etwas, das auf keinen Fall in der Reflexion stehen soll?\par
\par\noindent\emph{Antwort:} das irgendwas mit KI gemacht worde oder das KI irgendwas hiervon geschrieben hat. dementsprechend sollte auch der schreibstil meinem ähnelt, also nicht viel zu professionell und eher gelassen.\par\smallskip
\par\noindent \textbf{J3} Sonst noch was, das dir zum Arbeitsplan wichtig ist?\par
\par\noindent\emph{Antwort:} vergleiche noch mit dem von janosch (angehängt im chat), das bedeutet steigere das ''professionalitäts''level nicht viel weiter als das von janosch\par\smallskip
}
